\section{El Problema y Cambio de Paradigma}

% --- Ponente 1 ---

\begin{frame}{El límite del diseño humano (\textit{Handcrafted})}
\begin{alertblock}{Internet colapsa bajo el peso visual}
El 80\% del tráfico de Internet es contenido visual (imágenes y vídeo).
\end{alertblock}
\begin{itemize}
    \item Los códecs clásicos (JPEG, HEVC) se basan en la \textbf{DCT} y procesamiento por bloques.
    \item Diseñados manualmente hace décadas, han alcanzado su \textbf{techo matemático}.
    \item Resultado: \textbf{artefactos de bloque} visibles a tasas bajas de bits.
    \item No aprovechan la redundancia semántica de la imagen.
\end{itemize}
\end{frame}

\begin{frame}{El Paradigma del Aprendizaje (NIC)}
\begin{block}{Neural Image Compression}
Dejar que una \textbf{red neuronal profunda} aprenda automáticamente la representación más eficiente de la imagen, sin reglas diseñadas a mano.
\end{block}
\vspace{0.5cm}
\textbf{Pregunta clave:}
\begin{enumerate}
    \item ¿Es \textbf{técnicamente} viable sustituir a JPEG hoy?
    \item ¿Es \textbf{industrialmente} viable su despliegue a escala?
\end{enumerate}
\end{frame}
